\section{Introduction}
Brian Kernighan once wrote "Everyone knows that debugging is twice as hard as 
writing a program in the first place. So if you're as clever as you can be when
you write it, how will you ever debug it?" \cite{Kernighan82}. A few years
earlier, he also wrote "the most effective debugging tool is still careful
thought, coupled with judiciously placed print statements" \cite{Kernighan78}.
The world of computer programming has advanced enormously in the time since,
yet debugging remains one of the greatest challenges software engineers face
on a daily basis. All this despite a great deal of work over the decades
aimed at solving the problem.

Attempts to make programs easier to debug have taken many forms over the years.
Debugging research can be broadly divided into two camps. The first are those
who focus on building general tools which can make debugging easier, like the 
GNU debugger project, or modern IDEs. The second camp espouse the design of 
languages and stricter programming style in order to make programs easier to
reason about. Proponents of pure functional programming often fall into this camp,
with the maxim "well-typed programs don't go wrong". In this paper, we straddle the 
gap, presenting a graphical framework which can be applied in multiple settings, and a 
research language we have implemented to demonstrate the efficacy of the framework.

The rest of the paper is structured as follows. In \ref{sec:background}, we present
motivating examples and break down the workflow of the system. In \ref{sec:prelims},
we cover the mathematical preliminaries needed for our approach. In \ref{sec:baseline}
we present the semantics of the language before the introduction of the graphical
elements in \ref{sec:graph}.
\clearpage